Assume \(L>0\).  Write
\[
 c_{n,d}=n^{-1/2}d^{-1/4},
 \qquad
 g(h)=Lh^2+c_{n,d}h^{-3/2},
 \qquad
 \ell=d^{-1},
 \qquad
 u=d^{-1/2}.
\tag{10}
\]
In the logarithmic-degree asymptotic regime defining the model class, \(d\geq4\) for all sufficiently large \(n\).  Thus \({\rm def:bandwidth-set}\) gives \(\mathcal H_{n,d}=[\ell,u]\) for all such \(n\).  By \({\rm def:optimal-bandwidth}\),
\[
 h_0^7=\frac{L^{-2}}{n\sqrt d}.
\]
Since \(c_{n,d}^2=(n\sqrt d)^{-1}\), taking positive square roots gives the exact balance identity
\[
 Lh_0^2=c_{n,d}h_0^{-3/2}.
\tag{11}
\]

Suppose first that \(h_0\leq\ell\).  For every \(h\in[\ell,u]\), \(g(h)\geq Lh^2\geq L\ell^2\).  Moreover, monotonicity of \(h\mapsto h^{-3/2}\), (11), and \(h_0\leq\ell\) imply
\[
 c_{n,d}\ell^{-3/2}
 \leq c_{n,d}h_0^{-3/2}
 =Lh_0^2
 \leq L\ell^2.
\]
Evaluating at \(\ell\) therefore gives
\[
 Ld^{-2}=L\ell^2
 \leq\inf_{h\in[\ell,u]}g(h)
 \leq g(\ell)
 \leq2L\ell^2=2Ld^{-2}.
\tag{12}
\]

Suppose next that \(\ell<h_0<u\).  Equation (11) gives \(g(h_0)=2Lh_0^2\).  If \(h\geq h_0\), then \(Lh^2\geq Lh_0^2\); if \(h\leq h_0\), then \(c_{n,d}h^{-3/2}\geq c_{n,d}h_0^{-3/2}=Lh_0^2\).  Hence
\[
 Lh_0^2
 \leq\inf_{h\in[\ell,u]}g(h)
 \leq2Lh_0^2.
\tag{13}
\]
Direct substitution of the definition of \(h_0\) yields
\[
 Lh_0^2
 =L\left(\frac{L^{-2}}{n\sqrt d}\right)^{2/7}
 =L^{3/7}(n^2d)^{-1/7}.
\tag{14}
\]

Finally, suppose that \(h_0\geq u\).  For every \(h\in[\ell,u]\),
\[
 g(h)\geq c_{n,d}h^{-3/2}
 \geq c_{n,d}u^{-3/2}
 =\sqrt{d/n}.
\tag{15}
\]
Furthermore, (11) and \(u\leq h_0\) give
\[
 Lu^2\leq c_{n,d}u^{-3/2}.
\]
Evaluation at \(u\) and (15) therefore show
\[
 \sqrt{d/n}
 \leq\inf_{h\in[\ell,u]}g(h)
 \leq g(u)
 \leq2\sqrt{d/n}.
\tag{16}
\]
By \({\rm def:frontier-envelope}\), \(\mathfrak r^{\mathrm{sup}}_{n,d,L}=\sqrt{d/n}+\inf_{h\in[\ell,u]}g(h)\).  Adding \(\sqrt{d/n}\) to (12), (13) together with (14), and (16), respectively, proves the three asserted \(\asymp\) relations, with absolute comparison constants.

It remains to verify the decision-theoretic statements.  The explicit inequality proved in \({\rm thm:honest-upper}\) is
\[
 \sup_{P\in\mathcal M_{n,d,L}}
 \mathbb E_P|\widehat\theta_{h^\dagger}-\theta_n|
 \leq\mathfrak u_{n,d,L}.
\tag{17}
\]
The procedure \(\widehat\theta_{h^\dagger}\) is a single measurable data-based estimator by \({\rm def:estimator}\) and \({\rm def:exact-optimal-bandwidth}\).  It is therefore an admissible competitor in the estimator infimum in \({\rm def:minimax-criteria}\), so (17) gives
\[
 R^*_{n,d,L}\leq\mathfrak u_{n,d,L}.
\tag{18}
\]
Likewise, \({\rm thm:honest-upper}\) proves that \(\mathcal C_{h^\dagger}\) meets the uniform \(1-\alpha\) coverage constraint in \({\rm def:minimax-criteria}\), and its expected length obeys the explicit bound in (9).  It is consequently an admissible competitor in the confidence-set infimum, which proves
\[
 \Lambda^*_{n,d,L}(\alpha)
 \leq2\max\{1,\alpha^{-1/2}\}\mathfrak u_{n,d,L}.
\tag{19}
\]

Lastly, \({\rm thm:conditional-isolate-wall}\) applies to every graph-only conditional-bias certificate over the same unchanged model class and gives, on each realized graph,
\[
 b(A)\geq\frac{B}{n}
 \sum_{i\in\mathcal V_n}\mathbf1\{N_i=0\}.
\tag{20}
\]
Thus the asserted unavoidable isolate term follows directly.  Equations (18)--(20) provide only a minimax upper bound and a conditional-identification obstruction.  They invoke neither a supported-unit lower experiment nor any distributional information beyond the tail first-moment summary \(\pi_{m,n}\); accordingly they do not imply a matched converse or determine a nonlinear missing-sign scale.  This proves precisely the positive claims and the scope qualification in the proposition.